\section{Related work}\label{sec:related_work}

\noindent\textbf{Execution-based Coverage and Test Adequacy.}
Statement and branch coverage remain the dominant measures of test adequacy, yet their relationship with fault-detection effectiveness is known to be limited~\cite{papadakis2019mutation, soremekun2023evaluating,inozemtseva2014coverage,gopinath2014code}. Test suites with similar execution coverage may exhibit substantially different defect-detection capabilities, especially for UI component libraries where correctness depends on behavioral relations among component properties, interactions, and rendering outcomes. Our work complements execution-based coverage by assessing whether inferred behavioral relations are exercised and explicitly validated.

\vspace{2pt}
\noindent\textbf{Metamorphic Testing and MR Inference.}
Metamorphic Testing (MT) addresses the oracle problem by specifying expected relations between multiple executions rather than absolute outputs~\cite{chen2020metamorphic, segura2016survey, liu2013effectively}. A long-standing challenge is the construction of meaningful metamorphic relations~\cite{chen2018metamorphic}. Recent studies have shown that large language models can effectively infer MRs from source code, documentation, and natural-language artifacts~\cite{NolascoMDGGPUAF24,zhang2023automated,shin2024towards, ayerdi2024genmorph, duque2024selecting, le2025can, ma2025specgen}. Existing work primarily uses inferred MRs for follow-up test generation or oracle construction. In contrast, we use inferred MRs as a behavioral reference for assessing the adequacy of existing UI component test suites.

\vspace{2pt}
\noindent\textbf{GUI and Component Testing.}
Prior work on GUI testing has focused on model-based, event-driven, and automated exploration techniques, with particular attention to the oracle problem~\cite{memon2001hierarchical,xie2007designing,barr2014oracle,wang2020combodroid,wang2021vet,ahlgren2021testing}. Existing studies of UI component testing mainly evaluate test quality through structural metrics, test smells, or co-evolution analyses rather than behavioral specifications~\cite{spadini2018relation,rwemalika2023smells}. Our work instead introduces a relation-level adequacy perspective based on inferred MRs.

\vspace{2pt}
\noindent\textbf{Relation to Other Adequacy Techniques.}
While MR coverage builds on metamorphic relations, its objective differs from mutation testing, checked coverage, and property-based testing. Rather than generating tests, injecting faults, or measuring assertion reachability, it uses inferred behavioral relations as a reference for assessing the adequacy of existing UI component test suites. Thus, MR coverage provides a complementary relation-level perspective on behavioral validation.